\section{本讲定位：为什么 CS336 还要从零构建语言模型}

Lecture 1 是整门 CS336 的路线图。它没有直接进入 Transformer 公式，而是先回答一个更根本的问题：当 GPT、Claude、Gemini 这类 API 已经非常强时，为什么研究者还要自己实现 tokenizer、Transformer、optimizer、training loop、kernel、parallelism 和 inference stack？

课程给出的答案是：抽象层提高了生产力，但 LLM 的抽象层仍然非常漏。Prompt、tokenization、context length、sampling、latency、cost、failure mode、tool use、alignment 都会把底层机制重新暴露出来。只会调用 API 的研究者可以很快构建 demo，但很难回答系统为什么失败、如何扩展、如何调试、如何做真正新的研究。

\begin{figure}[H]
\centering
\includegraphics[width=0.9\textwidth]{course-staff.png}
\caption{CS336 Spring 2026 课程团队。本讲开场强调这是第三次开课，课程仍坚持 from-scratch philosophy。}
\figsource{CS336 Spring 2026 Lecture 1 官方可执行讲义，\texttt{welcome()}。}
\end{figure}

\begin{knowledgebox}{读图：课程团队页为什么重要}
这张图不是普通封面。它告诉读者：CS336 是一门持续迭代的研究工程课程。2026 版仍保留“从零构建”的主线，但增加更多现代 LM ingredients，如 mixture of experts、long-context、agents。也就是说，本课不是复古地重写一个小 Transformer，而是用可控规模理解现代大模型系统。
\end{knowledgebox}

\begin{importantbox}{第一讲的核心问题}
在固定资源下，如何训练出最好的语言模型？这里的资源不只是 GPU FLOPs，还包括数据、显存、通信带宽、工程时间、推理预算和可调试性。第一讲所有内容都围绕这个问题展开。
\end{importantbox}

\subsection{API 抽象为什么不够}

讲义把研究者和底层技术的距离变化压缩成三步：

\begin{longtable}{p{0.18\textwidth}p{0.34\textwidth}p{0.38\textwidth}}
\toprule
\textbf{阶段} & \textbf{常见工作方式} & \textbf{能力与局限} \\
\midrule
2016 & 自己实现并训练模型 & 能理解机制，但规模较小、工程成本高。 \\
2018 & 下载 BERT 等模型并 fine-tune & 生产力提升，但许多底层选择已经被预设。 \\
今天 & 调用 GPT/Claude/Gemini API & 应用构建更快，但很多 failure mode 无法解释。 \\
\bottomrule
\end{longtable}

\begin{warningbox}{只会使用 API 的局限}
API 使用者很难回答这些问题：为什么 tokenizer 改变会影响验证集 loss？为什么同样参数量的模型训练速度不同？为什么长上下文推理突然变慢？为什么某些 prompt 的失败无法用“模型不够聪明”解释？这些问题都要求理解底层 stack。
\end{warningbox}

\subsection{什么是 executable lecture}

CS336 2026 的 Percy 讲义经常以可执行 Python source 发布。`text(...)`、`image(...)`、`link(...)` 不是普通 slide 标记，而是构成 lecture trace 的节点；`@inspect`、`@stepover` 让学生看到代码执行中的变量状态。

\begin{lstlisting}[caption={executable lecture 的最小心智模型}]
def lecture():
    text("Explain an idea")
    image("diagram.png")
    value = compute()  # @inspect value
    text("Interpret the result")
\end{lstlisting}

\begin{knowledgebox}{为什么 executable lecture 适合 CS336}
传统 slide 展示结论，代码展示实现，讲者口头解释连接二者。Executable lecture 把这三件事放进同一个 trace：可以读结构、看图、跑代码、检查变量。这与“understanding via building”的课程哲学一致。
\end{knowledgebox}

\subsection{本章小结}

CS336 的第一讲先建立学习姿态：不要把 LLM 当远程 API，也不要把小模型当玩具。课程希望学生在可控规模下掌握可迁移的 mechanics 和 mindset，然后谨慎处理 scale-specific intuition。

